\section{Related Work}
\label{sec:related}

\paragraph{LLM-based topic discovery.}
TopicGPT~\citep{pham2024topicgpt} prompts an LLM to generate
human-readable topic labels and short definitions directly from sampled
documents, then assigns each remaining document to one of the generated
topics. It demonstrated that LLM-produced topics are substantially more
interpretable than classical BERTopic clusters, but the workflow is
essentially a two-stage prompt: there is no typed edit vocabulary, no
explicit refinement loop with validation, and no first-class notion of
cost or partial progress. TnT-LLM~\citep{wan2024tntllm} from Microsoft
stages a heavier two-phase pipeline in which an LLM first summarises
batches of a conversational corpus into candidate labels and a second
phase distils a lightweight classifier that scales the labelling to the
full corpus. TnT-LLM is closer in spirit to our setting -- it uses LLMs
as label generators and as labellers -- but it is structured as an
offline production pipeline rather than an interactive agent, the
``refine'' step is a single summarisation pass over batches rather than
an iterative critique loop driven by per-item don't-fit signals, and
the system is not released as an open tool with live cost accounting
or a UI.

\paragraph{Classical topic modelling.}
BERTopic~\citep{grootendorst2022bertopic} represents the strongest
non-LLM baseline: it embeds documents with a sentence transformer,
reduces dimensionality with UMAP, clusters with HDBSCAN, and labels
clusters via class-based TF-IDF. It is fast, fully reproducible, and
needs no API key, but its outputs are bag-of-keyword cluster
descriptions rather than human-readable categories with definitions,
and the taxonomy cannot be edited in a typed way once produced.
Classical LDA~\citep{blei2003lda} remains a reference point but
suffers from the same interpretability and editability issues at a
larger margin.

\paragraph{Goal-driven and human-in-the-loop systems.}
GoalEx~\citep{wang2023goalex} formulates explanation as goal-driven
clustering: a user specifies a natural-language goal and the system
produces a set of explanations partitioning the corpus that satisfy
that goal. LLooM~\citep{lam2024lloom} is a CHI~2024 mixed-initiative
system that surfaces high-level concepts from unstructured text for
analysts to inspect and edit interactively. Co-DETECT~\citep{codetect2025}
is an EMNLP~2025 demonstration that supports collaborative codebook
development with LLMs, with the human and the model alternating
proposals and revisions. These systems share our concern with
interactive refinement, but they target the human-loop side of the
problem; \textsc{taxonomy\_agent} is complementary in that it
formalises the \emph{model-side} edit loop -- including a typed revise
vocabulary and a self-reported convergence signal -- which a future
human-in-the-loop layer can sit on top of.

\paragraph{Recent 2025 work.}
TaxoAdapt~\citep{taxoadapt2025} adapts an existing taxonomy to a new
domain via LLM-guided edits, sharing our intuition that taxonomy
construction is naturally an edit problem rather than a generation
problem; it focuses on taxonomy \emph{adaptation} given a seed,
whereas we discover from scratch under a typed DSL.
LLM-guided semantic-aware clustering~\citep{llmguided2025} uses an LLM
to steer the clustering objective itself, blurring the line between
classical clustering and LLM topic discovery. The recent EMNLP~2024
survey on LLMs for data annotation~\citep{tan2024llmannot} places this
family of work in the broader annotation landscape and motivates the
operational concerns -- cost, reproducibility, interruption-safety --
that our system targets.

\paragraph{Positioning of \textsc{taxonomy\_agent}.}
Against TopicGPT and the single-shot LLM baseline, our system adds an
explicit iterative loop with a typed edit vocabulary and a
convergence criterion, rather than a one-shot prompt. Against TnT-LLM,
we expose the refinement step as a small set of validated tool calls
driven by per-item don't-fit signals, separate the expensive
orchestrator from a cheap parallel judge, and ship the system as an
interactive open-source demo with live cost accounting and
partial-progress saving rather than an offline production pipeline.
Against BERTopic and LDA, we produce human-readable categories with
definitions and support post-hoc typed edits. Against goal-driven and
human-in-the-loop systems (GoalEx, LLooM, Co-DETECT), we focus on the
model-side edit loop and provide infrastructure -- the six-op DSL with
validate-before-mutate, the orchestrator--judge split, the cost and
progress streams -- on which interactive human workflows can be
layered. Against TaxoAdapt and LLM-guided clustering, we operate
without a seed taxonomy and externally validate our self-reported
convergence signal. The combination -- orchestrator--judge separation,
typed validated revise DSL, externally-validated self-reported
convergence, native cost tracking, and an open-source model-agnostic
Streamlit demo -- is, to our knowledge, novel.
